\section{演示数据的采集与质量保障}

\subsection{数据采集管线}

演示数据来自人类、仿真机器人或 pre-recorded log。常见 pipeline：
\begin{itemize}
    \item \textbf{Kinesthetic Teaching}：人体直接控制机械臂，低 latency。
    \item \textbf{Remote teleoperation}：网格化控制，支持地理分布式专家。
    \item \textbf{Behavioral cloning from screen recording}：视频+键盘鼠标行为同步。
\end{itemize}

\begin{knowledgebox}{embodiment gap}
直接用人类视频代替 robot demonstration 会产生 embodiment gap（自由度差、动力学差），但可以作为 exploration signal，通过 domain randomization 缩小 gap。
\end{knowledgebox}

\subsection{数据 QA 与 audit}

高质量演示需要 audit pipeline，包括：
\begin{itemize}
    \item 专家 agreement rate（Cohen's κ）；
    \item 状态-动作对的 variance coverage；
    \item 误差示范的 labeling；
    \item Metadata（任务难度、环境 lighting）。
\end{itemize}

\begin{table}[H]
\centering
\begin{tabular}{p{0.28\textwidth} p{0.35\textwidth} p{0.33\textwidth}}
\toprule
维度 & 指标 & 工程应对 \\
\midrule
一致性 & Pairwise agreement & 不一致样本送审 \\
多样性 & Different skills, contexts & 用 clustering 识别覆盖 gaps \\
安全性 & Adversarial prompt & 加入 counterfactual data \\
\bottomrule
\end{tabular}
\caption{演示数据质量保障}
\end{table}

\subsection{跨模态与跨角色协作}

现代模仿学习项目往往涉及数据标注团队、语言专家与 SRE，共同维护 dataset catalog 与 manifest。每次 rollout 都需要把 new prompts 加入 dataset，并把 shift log 记录在 evidence board 中。

\subsection{本章小结}

演示数据质量直接决定模仿学习的下限。通过严格 audit、metadata 和 multi-role feedback，可以把 embodiment gap、adversarial prompt 等问题转化为可治理的风险。

